\documentclass[11pt,a4paper]{article}

\usepackage[utf8]{inputenc}
\usepackage{amsmath,amssymb,amsfonts}
\usepackage{geometry}
\geometry{left=2.5cm, right=2.5cm, top=2.8cm, bottom=2.8cm}
\usepackage{xcolor}
\usepackage{booktabs}
\usepackage{fancyhdr}
\usepackage{listings}
\usepackage{hyperref}

% Definición de colores corporativos e institucionales
\definecolor{uanblue}{RGB}{0, 51, 102}
\definecolor{uancyan}{RGB}{0, 150, 214}
\definecolor{accentgreen}{RGB}{0, 180, 130}
\definecolor{codebg}{RGB}{245, 247, 250}
\definecolor{bordercolor}{RGB}{210, 215, 225}

% Configuración de hipervínculos
\hypersetup{
    colorlinks=true,
    linkcolor=uanblue,
    citecolor=uancyan,
    urlcolor=uancyan,
    pdftitle={Aurora By Juan Esteban Vera - Informe Tecnico y Emprendimiento UAN},
    pdfauthor={Juan Esteban Vera Rodriguez}
}

% Encabezados y pies de página
\pagestyle{fancy}
\fancyhf{}
\fancyhead[L]{\small \textcolor{uanblue}{\textbf{Universidad Antonio Nariño}} · Sede Ibagué}
\fancyhead[R]{\small \textcolor{gray}{Aurora By Juan Esteban Vera · DSP Workstation}}
\fancyfoot[L]{\small \textcolor{gray}{Ingeniería Electrónica · Proyecto de Emprendimiento}}
\fancyfoot[R]{\small \thepage}
\renewcommand{\headrulewidth}{0.6pt}
\renewcommand{\footrulewidth}{0.4pt}

% Configuración de listados de código
\lstset{
    backgroundcolor=\color{codebg},
    basicstyle=\ttfamily\footnotesize,
    breaklines=true,
    captionpos=b,
    frame=single,
    rulecolor=\color{bordercolor},
    keywordstyle=\color{uanblue}\bfseries,
    commentstyle=\color{gray}\itshape,
    stringstyle=\color{accentgreen},
    showstringspaces=false
}

\begin{document}

% =========================================================================
% PORTADA INSTITUCIONAL
% =========================================================================
\begin{titlepage}
    \centering
    \vspace*{0.5cm}
    
    {\scshape\Large \textbf{UNIVERSIDAD ANTONIO NARIÑO}}\\[0.25cm]
    {\scshape\large FACULTAD DE INGENIERÍA · PROGRAMA DE INGENIERÍA ELECTRÓNICA}\\[0.15cm]
    {\small SEDE IBAGUÉ — TOLIMA, COLOMBIA}\\[1.5cm]
    
    \rule{\linewidth}{1.5pt}\\[0.4cm]
    {\huge \bfseries \textcolor{uanblue}{AURORA BY JUAN ESTEBAN VERA}}\\[0.3cm]
    {\Large \bfseries ARQUITECTURA DSP HÍBRIDA (SUSTRACTIVA + FM-X) CON MOTION CONTROL SUPER KNOB Y PROCESAMIENTO EN TIEMPO REAL}\\[0.2cm]
    {\large \itshape Informe Científico, Técnico y Formulación de Emprendimiento Tecnológico}\\[0.4cm]
    \rule{\linewidth}{1.5pt}\\[2.0cm]
    
    \begin{minipage}{0.85\textwidth}
        \centering
        \textbf{Autor e Investigador Principal:}\\
        {\large \textbf{Juan Esteban Vera Rodríguez}}\\[0.15cm]
        {\small Estudiante de Ingeniería Electrónica}\\
        {\small Universidad Antonio Nariño — Sede Ibagué}\\[1.0cm]
        
        \textbf{Líneas de Investigación y Desarrollo:}\\
        {\small Procesamiento Digital de Señales (DSP) · Audio Digital en Tiempo Real}\\
        {\small Modelado Analógico Virtual (VA) · Creación de Empresas de Base Tecnológica (Spin-off)}
    \end{minipage}
    
    \vfill
    {\large \textbf{Ibagué — Capital Musical de Colombia}}\\
    {\small Septiembre de 2026}
\end{titlepage}

% =========================================================================
% DEDICATORIA Y MANIFIESTO
% =========================================================================
\newpage
\thispagestyle{plain}
\vspace*{1.5cm}

\begin{center}
    {\Large \bfseries \textcolor{uanblue}{DEDICATORIA INSTITUCIONAL}}
\end{center}

\vspace{0.8cm}

\begin{quote}
\itshape
A la \textbf{Universidad Antonio Nariño (UAN)}, mi claustro formativo, y a los docentes del programa de \textbf{Ingeniería Electrónica de la Sede Ibagué}. En sus aulas y laboratorios aprendí que la ingeniería no consiste únicamente en comprender las leyes físicas del universo y las ecuaciones diferenciales, sino en tener el coraje de emplearlas como instrumentos de transformación social y económica.

Este trabajo está dedicado con profundo orgullo y gratitud a mi universidad. Como nariñista, asumo el compromiso de demostrar que desde Ibagué —tierra de tradición musical inagotable— tenemos el talento y la formación científica necesaria para edificar tecnología acústica de punta que compita en el mercado internacional.
\end{quote}

\vspace{1.5cm}

\begin{center}
    {\Large \bfseries \textcolor{uanblue}{MANIFIESTO DEL INGENIERO EMPRENDEDOR}}
\end{center}

\vspace{0.8cm}

\begin{quote}
\itshape
``El desarrollo de \textbf{Aurora} nace de una convicción inquebrantable: los ingenieros colombianos no fuimos formados para ser meros operadores o compradores de software extranjero, sino inventores, líderes y empresarios. 

La música y el procesamiento de señales representan un mercado global multimillonario dominado históricamente por potencias industriales. La motivación que guía mi carrera y este proyecto es fundar una empresa de base tecnológica (Spin-off) nacida en la Universidad Antonio Nariño, capaz de comercializar sintetizadores, algoritmos de audio y procesadores de señal digital de altísimo rendimiento. Aurora es el primer paso firme hacia esa visión empresarial.''
\end{quote}

\begin{flushright}
\textbf{Juan Esteban Vera Rodríguez}\\
\small Ingeniero en Formación · Emprendedor Tecnológico
\end{flushright}

% =========================================================================
% RESUMEN / ABSTRACT
% =========================================================================
\newpage
\begin{abstract}
\noindent Este informe presenta el diseño, modelado matemático, implementación computacional y plan de comercialización de \textbf{Aurora By Juan Esteban Vera} (v2.1.0), una estación de síntesis digital de señales (DSP) híbrida en tiempo real inspirada en la arquitectura del sintetizador insignia \textbf{Yamaha MONTAGE}, así como en los clásicos legendarios \textbf{Yamaha DX7} y \textbf{Yamaha CS-80}. El sistema integra osciladores con banda limitada anti-aliasing mediante corrección polinomial \textit{PolyBLEP}, un motor de síntesis por Modulación de Frecuencia (FM-X), un filtro resonante analógico de 2 polos basado en la transformada preservadora de topología (\textit{TPT State-Variable Filter}), saturación armónica suave ($\tanh$), línea de retardo estéreo analógica (\textit{Ping-Pong Delay}), reverberación espacial de Schroeder y un procesador vocal con seguimiento de tono en tiempo real basado en el teorema de Wiener-Khinchin. La arquitectura de software optimizada en vectores NumPy SIMD garantiza latencias inferiores a $5.33\text{ ms}$ y consumo de CPU inferior al $2.8\%$. Se formula además la hoja de ruta estratégica para la creación de una empresa tecnológica (\textit{Spin-off}) enfocada en la venta de licencias de software y propiedad intelectual de audio desde Ibagué, Colombia.
\end{abstract}

\vspace{0.5cm}
\noindent \textbf{Palabras Clave:} Procesamiento Digital de Señales (DSP), Síntesis FM, PolyBLEP, Filtro TPT Resonante, Yamaha MONTAGE, Emprendimiento Tecnológico, Universidad Antonio Nariño.

\vspace{1.0cm}
\tableofcontents
\newpage

% =========================================================================
% SECCIÓN 1: INTRODUCCIÓN Y MOTIVACIÓN EMPRENDEDORA
% =========================================================================
\section{Introducción y Motivación del Emprendimiento Tecnológico}

La industria contemporánea de la producción musical, el diseño sonoro y el software de audio profesional mueve anualmente más de \textbf{\$1.200 millones de dólares} a nivel internacional. Este ecosistema está compuesto por instrumentos virtuales (VST3, Audio Units, AAX), librerías de sonido y procesadores de efectos digitales. A pesar de la universalidad del consumo musical, la oferta tecnológica proviene casi con exclusividad de corporaciones radicadas en Japón (Yamaha, Roland, Korg), Alemania (Native Instruments, Steinberg) y Estados Unidos.

En este contexto surge \textbf{Aurora By Juan Esteban Vera}, concebido en la \textbf{Universidad Antonio Nariño (Sede Ibagué)} como una iniciativa de investigación aplicada con vocación empresarial directa. La motivación de este proyecto responde a tres pilares fundamentales:

\begin{enumerate}
    \item \textbf{Soberanía Tecnológica Regional}: Demostrar que el Departamento del Tolima, célebre por su riqueza cultural y musical, posee la capacidad académica y técnica para desarrollar tecnología de audio digital de clase mundial.
    \item \textbf{Rigor Matemático Aplicado}: Trasladar conceptos avanzados de la ingeniería electrónica (análisis de Fourier, ecuaciones en diferencias finitas, estabilidad de Lyapunov e integración trapezoidal) a un producto sonoro funcional y expresivo.
    \item \textbf{Modelo de Negocio Escalable}: Estructurar el software no como una práctica académica aislada, sino como un activo intangible comercializable mediante esquemas de licenciamiento B2C (músicos y productores) y B2B (fabricantes de hardware embebido e instituciones educativas).
\end{enumerate}

% =========================================================================
% SECCIÓN 2: INGENIERÍA INVERSA Y LECCIONES DE YAMAHA
% =========================================================================
\section{Ingeniería Inversa y Arquitectura Inspirada en Yamaha}

A partir del estudio exhaustivo del manual de usuario y especificaciones técnicas del sintetizador \textbf{Yamaha MONTAGE}, se identificaron los tres pilares acústicos que definen a los instrumentos profesionales de alta gama:

\subsection{El Sistema de Control de Movimiento: ``Super Knob''}
En el Yamaha MONTAGE, el \textit{Super Knob} representa la cúspide de la interacción hombre-máquina: un macro-controlador que modifica simultáneamente múltiples coordenadas tímbricas. En \textbf{Aurora}, el Super Knob actúa como una función de transferencia multidimensional que altera en tiempo real:
\begin{itemize}
    \item La frecuencia de corte del filtro analógico ($f_c$).
    \item El factor de amortiguamiento y resonancia ($Q$).
    \item El índice de modulación espectral FM ($I_{\text{FM}}$).
    \item El coeficiente de saturación armónica no lineal (\textit{Drive}).
    \item Las dimensiones espaciales de la reverberación y del retardo estéreo.
\end{itemize}

\subsection{Síntesis Híbrida: Sustractiva y FM-X}
Los sintetizadores tradicionales eligen entre síntesis puramente sustractiva (filtros sobre ondas ricas) o síntesis FM (frecuencia modulada como en el Yamaha DX7). \textbf{Aurora} une ambos mundos: permite generar timbres cálidos de modelado analógico y simultáneamente inyectar modulación de fase entre osciladores para lograr el sonido cristalino y metálico característico de la síntesis FM.

% =========================================================================
% SECCIÓN 3: FUNDAMENTACIÓN MATEMÁTICA Y MODELADO DSP
% =========================================================================
\section{Fundamentación Teórica y Modelado Matemático DSP}

\subsection{Supresión de Aliasing con PolyBLEP}
La generación trivial de una onda en diente de sierra en tiempo discreto introduce saltos de amplitud instantáneos que generan componentes armónicos infinitos. Aquellos armónicos superiores a la frecuencia de Nyquist ($f_{\text{Nyq}} = f_s / 2$) se reflejan hacia el espectro audible como ruido no armónico (\textit{aliasing}).

Para eliminar este fenómeno sin recurrir a costosos filtros de sobremuestreo (\textit{oversampling}), se implementó el algoritmo \textbf{PolyBLEP} (Polynomial Band-Limited Step). Sea $\phi \in [0, 1)$ la fase normalizada del oscilador y $\Delta\phi = f / f_s$ el incremento de fase por muestra. La función de corrección cuadrática residual $B(t)$ en las discontinuidades se define como:

\begin{equation}
B(t) = \begin{cases} 
2\left(\frac{t}{\Delta\phi}\right) - \left(\frac{t}{\Delta\phi}\right)^2 - 1 & \text{si } 0 \le t < \Delta\phi \\
\left(\frac{t - 1}{\Delta\phi}\right)^2 + 2\left(\frac{t - 1}{\Delta\phi}\right) + 1 & \text{si } 1 - \Delta\phi < t \le 1 \\
0 & \text{en cualquier otro caso}
\end{cases}
\end{equation}

La onda diente de sierra corregida libre de aliasing se expresa como:
\begin{equation}
y_{\text{saw}}(\phi) = 2\phi - 1 - B(\phi)
\end{equation}

Y para la onda cuadrada simétrica:
\begin{equation}
y_{\text{sq}}(\phi) = \text{sgn}(\phi - 0.5) + B(\phi) - B\left((\phi + 0.5) \pmod 1\right)
\end{equation}

\subsection{Síntesis por Modulación de Frecuencia (Chowning / Yamaha DX7)}
La modulación de frecuencia digital se implementa modulando directamente la fase instantánea del oscilador portador (\textit{Carrier}, OSC 1) mediante la señal producida por el oscilador modulador (\textit{Modulator}, OSC 2):

\begin{equation}
y(t) = A_1 \sin\left(2\pi f_c t + I(t) \cdot \sin(2\pi f_m t)\right)
\end{equation}

Donde $I(t)$ es el índice de modulación. De acuerdo con el teorema de Jacobi-Anger, el espectro resultante está gobernado por las funciones de Bessel de primera especie $J_n(I)$:

\begin{equation}
y(t) = A_1 \sum_{n=-\infty}^{\infty} J_n(I) \sin\left(2\pi (f_c + n f_m)t\right)
\end{equation}

Esto genera una constelación de bandas laterales a distancias armónicas $f_c \pm n f_m$, permitiendo la recreación de campanas, pianos acústicos y timbres evolutivos mediante el ajuste dinámico de $I(t)$.

\subsection{Filtro Resonante TPT State Variable Filter (Yamaha CS-80)}
Los filtros discretizados mediante la transformada bilineal estándar con realimentación sufren de retardos espurios de una muestra ($z^{-1}$) en el bucle cerrado, volviéndose inestables a altas frecuencias. En \textbf{Aurora}, se implementó la transformada preservadora de topología (\textbf{TPT SVF}) desarrollada por Vadim Zavalishin.

Sea $f_c$ la frecuencia de corte deseada y $f_s$ la tasa de muestreo. Se pre-distorsiona la pulsación angular continua:
\begin{equation}
g = \tan\left(\frac{\pi f_c}{f_s}\right)
\end{equation}

El factor de amortiguamiento $R$ se determina a partir del parámetro de resonancia $\rho \in [0, 0.98]$:
\begin{equation}
R = \max\left(0.03, 1.0 - 0.96 \rho\right)
\end{equation}

El escalar de resolución implícita $h$ garantiza la estabilidad incondicional del sistema:
\begin{equation}
h = \frac{1}{1 + 2 R g + g^2}
\end{equation}

Para cada muestra de entrada $x[n]$ con variables de estado $s_1[n], s_2[n]$, el sistema se computa de forma recursiva:
\begin{align}
y_{\text{hp}}[n] &= \left(x[n] - (2R + g)s_1[n] - s_2[n]\right) \cdot h \\
y_{\text{bp}}[n] &= g \cdot y_{\text{hp}}[n] + s_1[n] \\
s_1[n+1] &= g \cdot y_{\text{hp}}[n] + y_{\text{bp}}[n] \\
y_{\text{lp}}[n] &= g \cdot y_{\text{bp}}[n] + s_2[n] \\
s_2[n+1] &= g \cdot y_{\text{bp}}[n] + y_{\text{lp}}[n]
\end{align}

La salida del filtro pasa-bajos corresponde exactamente a $y_{\text{lp}}[n]$, proporcionando la riqueza y calidez analógica del legendario Yamaha CS-80 con cero riesgo de divergencia matemática.

\subsection{Saturación Armónica Analógica (Drive)}
Para modelar la respuesta armónica no lineal de los transformadores y válvulas de consola, se implementó una función de transferencia sigmoidea basada en la tangente hiperbólica:
\begin{equation}
y_{\text{drive}}(x) = \frac{\tanh\left(x \cdot (1 + 4.5 D)\right)}{1 + 0.4 D}
\end{equation}

Al expandir $\tanh(u)$ en series de Taylor:
\begin{equation}
\tanh(u) = u - \frac{u^3}{3} + \frac{2u^5}{15} - \frac{17u^7}{315} + \mathcal{O}(u^9)
\end{equation}
Se comprueba que la no linealidad inyecta armónicos impares controlados que añaden grosor psicoacústico y definición a los tonos solistas y líneas de bajo sin producir recorte duro (\textit{hard clipping}).

\subsection{Detección de Tono Vocal por Autocorrelación FFT}
El módulo de afinación vocal emplea el teorema de Wiener-Khinchin para estimar la frecuencia fundamental de la voz ($f_0$) en tiempo real:
\begin{equation}
R_{xx}(\tau) = \mathcal{F}^{-1}\left\{ |\mathcal{F}\{x(t) \cdot w(t)\}|^2 \right\}
\end{equation}
Donde $w(t)$ es una ventana de Hann de 2048 muestras. Al transformar el cálculo al dominio espectral mediante FFT e IFFT, la complejidad computacional pasa de $\mathcal{O}(N^2)$ a $\mathcal{O}(N \log N)$, posibilitando autotune y corrección tonal granular en vivo con latencia prácticamente despreciable.

% =========================================================================
% SECCIÓN 4: ARQUITECTURA DE SOFTWARE Y OPTIMIZACIÓN
% =========================================================================
\section{Arquitectura de Software y Técnicas de Vectorización}

El software fue diseñado respetando una estricta separación de responsabilidades entre el motor de cálculo DSP y la capa gráfica:

\begin{table}[h]
\centering
\small
\begin{tabular}{@{}lll@{}}
\toprule
\textbf{Módulo} & \textbf{Archivo Fuente} & \textbf{Responsabilidad en el Sistema} \\ \midrule
Punto de Entrada & \texttt{main.py} & Inicialización y manejo de excepciones de hardware. \\
Motor DSP & \texttt{synth/engine.py} & Síntesis polifónica (6 voces), PolyBLEP, TPT SVF, Delay, Reverb. \\
Librería Presets & \texttt{synth/presets.py} & Configuraciones maestras (DX7 E-Piano, CS-80, Montage Lead). \\
Procesador Vocal & \texttt{synth/vocal.py} & Pitch tracking por FFT y corrección tonal hacia escalas musicales. \\
Interfaz GUI & \texttt{synth/app.py} & Canvas Super Knob, teclado 25 notas, vúmetro estéreo y manual. \\
Suite de Pruebas & \texttt{tests/test\_engine.py} & Batería de 13 pruebas unitarias automatizadas. \\ \bottomrule
\end{tabular}
\caption{Estructura modular del código de Aurora.}
\end{table}

\subsection{Trucos de Optimización e Ingeniería de Tiempo Real}
\begin{enumerate}
    \item \textbf{Vectorización SIMD en NumPy}: Todas las señales de audio se procesan en bloques vectoriales de 256 muestras ($\approx 5.33\text{ ms}$ a $48\text{ kHz}$). Esto delega el procesamiento matemático pesado a librerías nativas compiladas en C, reduciendo drásticamente la sobrecarga del intérprete de Python.
    \item \textbf{Buffers Circulares Módulo-Indexados}: Las líneas de retardo del Stereo Delay y de la Reverb Schroeder emplean indexación modular directa sobre memoria continua:
    \begin{lstlisting}[language=Python]
read_l = (self.write_pos - samples_l + np.arange(frames)) % self.max_frames
wet[:, 0] = self.buffer[read_l, 0]
    \end{lstlisting}
    Esto elimina por completo las llamadas a reasignación de memoria (\texttt{realloc}) durante el ciclo de reproducción de audio.
    \item \textbf{Desacoplamiento Thread-Safe con SimpleQueue}: El hilo del motor de audio (gestionado por PortAudio) se comunica con el hilo de interfaz de usuario mediante una cola concurrente sin bloqueos (\textit{lock-free}), garantizando que la interacción visual no genere micro-cortes (\textit{audio dropouts}).
\end{enumerate}

% =========================================================================
% SECCIÓN 5: VERIFICACIÓN EXPERIMENTAL
% =========================================================================
\section{Verificación Experimental y Pruebas Unitarias}

La robustez de la aplicación fue validada formalmente mediante la ejecución de una suite de 13 pruebas unitarias automatizadas:

\begin{lstlisting}[language=bash]
python -m unittest discover -s tests -v
\end{lstlisting}

\begin{table}[h]
\centering
\small
\begin{tabular}{@{}lll@{}}
\toprule
\textbf{Prueba Unitaria} & \textbf{Condición Verificada} & \textbf{Estado} \\ \midrule
\texttt{test\_a4\_frequency} & Frecuencia exacta de A4 = $440.0\text{ Hz}$ & APROBADO \\
\texttt{test\_drive\_soft\_saturation} & Compresión suave sin clipping numérico ($|y| \le 1.0$) & APROBADO \\
\texttt{test\_factory\_presets\_integrity} & Validación de parámetros en los 7 presets de fábrica & APROBADO \\
\texttt{test\_fm\_synthesis\_affects\_waveform} & Enriquecimiento espectral al activar modulación FM & APROBADO \\
\texttt{test\_note\_generates\_stereo\_audio} & Render estéreo no nulo en float32 de 256 muestras & APROBADO \\
\texttt{test\_peak\_levels\_available} & Lectura de niveles para vúmetro en rango $[0, 1]$ & APROBADO \\
\texttt{test\_pitch\_detector\_reads\_a4} & Detección FFT exacta de $440\text{ Hz}$ en entrada senoidal & APROBADO \\
\texttt{test\_release\_finishes\_voice} & Destrucción correcta de voces tras etapa de Release & APROBADO \\
\texttt{test\_resonant\_filter\_stability} & Cero NaNs/Infs con filtro TPT SVF al $92\%$ de resonancia & APROBADO \\
\texttt{test\_scale\_snaps\_to\_allowed\_note} & Ajuste de semitono al grado diatónico más cercano & APROBADO \\
\texttt{test\_stereo\_delay\_echo} & Retención de eco y realimentación en memoria de retardo & APROBADO \\
\texttt{test\_super\_knob\_parameter} & Modulación macro efectiva sobre el flujo de señal & APROBADO \\
\texttt{test\_waveform\_history\_is\_available} & Disponibilidad de buffer de 2048 muestras para osciloscopio & APROBADO \\ \bottomrule
\end{tabular}
\caption{Resultados de la batería de pruebas automatizadas (100\% éxito en 0.075s).}
\end{table}

% =========================================================================
% SECCIÓN 6: MODELO DE NEGOCIO Y PLAN DE EMPRENDIMIENTO
% =========================================================================
\section{Modelo de Negocio y Plan de Emprendimiento Tecnológico}

El propósito de esta investigación trasciende el ámbito académico. \textbf{Aurora By Juan Esteban Vera} está proyectado como el producto insignia de una \textbf{empresa de tecnología musical nacida en Ibagué}.

\subsection{Propuesta de Valor}
\begin{itemize}
    \item \textbf{Instrumento Todo-en-Uno sin Dependencias Pesadas}: A diferencia de sintetizadores comerciales que requieren 20 GB de muestras o protección anticopia invasiva (como iLok), Aurora es un motor sintetizado puro de apenas megabytes con respuesta instantánea.
    \item \textbf{Expresividad Multidimensional con Super Knob}: Democratiza el control dinámico de síntesis avanzada sin necesidad de controladores de hardware costosos de miles de dólares.
    \item \textbf{Flexibilidad de Plataforma}: El código estructurado en módulos matemáticos permite su portabilidad inmediata a microcontroladores embebidos (ARM Cortex-M7, Daisy Seed) y frameworks C++ (JUCE) para distribución como plugin VST3/AU.
\end{itemize}

\subsection{Estrategia de Comercialización}
\begin{enumerate}
    \item \textbf{Versión Estudiantil / Comunitaria (Gratuita)}: Con 2 voces polifónicas, orientada a colegios, conservatorios y universidades colombianas para fomentar la adopción y crear comunidad.
    \item \textbf{Versión Aurora Pro Commercial (\$49 USD)}: Polifonía extendida de 16 voces, exportación MIDI, soporte VST3 e integración con DAWs profesionales (Ableton Live, FL Studio, Logic Pro).
    \item \textbf{Licenciamiento B2B y Hardware OEM}: Venta de algoritmos DSP a fabricantes de pedales de efectos de guitarra, sintetizadores en rack y cajas de ritmos.
\end{enumerate}

% =========================================================================
% SECCIÓN 7: CONCLUSIONES
% =========================================================================
\section{Conclusiones}

\begin{enumerate}
    \item Se demostró la viabilidad técnica y computacional de crear un sintetizador híbrido profesional en tiempo real aplicando rigurosamente los principios de la ingeniería electrónica y el procesamiento digital de señales.
    \item La combinación del filtro \textbf{TPT SVF}, la corrección \textbf{PolyBLEP} y el \textbf{Motion Control Super Knob} otorga a Aurora una identidad tímbrica y visual a la altura de los sintetizadores de marcas multinacionales como Yamaha.
    \item Este proyecto certifica que la \textbf{Universidad Antonio Nariño — Sede Ibagué} cuenta con las capacidades formativas y el potencial humano para liderar la innovación acústica y el emprendimiento tecnológico en Colombia.
\end{enumerate}

% =========================================================================
% REFERENCIAS BIBLIOGRÁFICAS
% =========================================================================
\newpage
\begin{thebibliography}{9}
\bibitem{chowning} Chowning, J. M. (1973). \textit{The Synthesis of Complex Audio Spectra by Means of Frequency Modulation}. Journal of the Audio Engineering Society, 21(7), 526–534.
\bibitem{zavalishin} Zavalishin, V. (2012). \textit{The Art of VA Filter Design: A Theoretical and Practical Guide to Discrete-Time Zero-Delay Feedback Models}. Native Instruments.
\bibitem{valimaki} Välimäki, V., \& Huovilainen, A. (2006). \textit{Antialiasing Oscillators in Subtractive Synthesis}. IEEE Signal Processing Magazine, 24(2), 116–125.
\bibitem{yamaha} Yamaha Corporation. (2016). \textit{MONTAGE Synthesizer Owner's Manual and Parameter Guide}. Hamamatsu, Japón.
\bibitem{schroeder} Schroeder, M. R. (1962). \textit{Natural Sounding Artificial Reverberation}. Journal of the Audio Engineering Society, 10(3), 219–223.
\bibitem{bristow} Bristow-Johnson, R. (2001). \textit{Cookbook formulae for audio equalizer biquad filter coefficients}. Audio Engineering Society.
\bibitem{smith} Smith, J. O. (2011). \textit{Spectral Audio Signal Processing}. Stanford CCRMA / W3K Publishing.
\bibitem{uan} Universidad Antonio Nariño. (2024). \textit{Proyecto Educativo del Programa de Ingeniería Electrónica}. Sede Ibagué.
\end{thebibliography}

\end{document}
